% =====================================================================
%  Projet Final - Fondamentaux de Big Data (M126)
%  ENSA Tetouan - Departement IA & Digitalisation - 2025-2026
%  Sujet 6 : Analyse predictive de la charge reseau et optimisation
%            de la consommation energetique en temps reel
%
%  Sections : 15. Ethique & Conformite
%             17. Recommandations pour les professionnels
%             18. Perspectives d'amelioration
%             19. Conclusion
%
%  Compilation recommandee :  pdflatex main.tex (x2 pour les TOC/refs)
% =====================================================================

\documentclass[11pt,a4paper]{article}

% --- Encodage et langue --------------------------------------------------
\usepackage[utf8]{inputenc}
\usepackage[T1]{fontenc}
\usepackage{lmodern}
\usepackage[french]{babel}
\usepackage{csquotes}
\usepackage{microtype}

% --- Mise en page --------------------------------------------------------
\usepackage[margin=2.5cm,headheight=14pt]{geometry}
\usepackage{setspace}
\onehalfspacing
\usepackage{parskip}

% --- Couleurs (assorties au document du sujet) ---------------------------
\usepackage{xcolor}
\definecolor{titlered}{RGB}{210, 50, 50}
\definecolor{sectiongreen}{RGB}{0, 130, 90}
\definecolor{accentorange}{RGB}{230, 130, 30}
\definecolor{lightgray}{RGB}{240, 240, 240}
\definecolor{softblue}{RGB}{30, 80, 140}
\definecolor{tableheaderbg}{RGB}{0, 130, 90}

% --- Typographie des titres ----------------------------------------------
\usepackage{titlesec}
\titleformat{\section}
  {\Large\bfseries\color{sectiongreen}}{\thesection}{0.6em}{}
\titleformat{\subsection}
  {\large\bfseries\color{accentorange}}{\thesubsection}{0.5em}{}
\titleformat{\subsubsection}
  {\normalsize\bfseries\color{softblue}}{\thesubsubsection}{0.4em}{}

% --- Tableaux ------------------------------------------------------------
\usepackage{booktabs}
\usepackage{tabularx}
\usepackage{longtable}
\usepackage{array}
\usepackage{multirow}
\usepackage{colortbl}
\renewcommand{\arraystretch}{1.3}
\newcolumntype{Y}{>{\raggedright\arraybackslash}X}
\newcolumntype{C}{>{\centering\arraybackslash}X}
\newcolumntype{L}[1]{>{\raggedright\arraybackslash}p{#1}}

% --- Listes --------------------------------------------------------------
\usepackage{enumitem}
\setlist{itemsep=2pt, parsep=2pt, topsep=3pt}

% --- Symboles ------------------------------------------------------------
\usepackage{pifont}
\newcommand{\done}{\textcolor{sectiongreen}{\ding{51}}}
\newcommand{\todo}{\textcolor{titlered}{\ding{55}}}

% --- Hyperliens ----------------------------------------------------------
\usepackage{hyperref}
\hypersetup{
  colorlinks=true,
  linkcolor=sectiongreen,
  urlcolor=accentorange,
  citecolor=sectiongreen,
  pdftitle={Projet Big Data - ENSA Tetouan},
  pdfauthor={Groupe Big Data}
}

% --- En-tete et pied de page --------------------------------------------
\usepackage{fancyhdr}
\pagestyle{fancy}
\fancyhf{}
\fancyhead[L]{\small\textbf{ENSA-Tétouan} -- Dépt IA \& Digitalisation}
\fancyhead[R]{\small Fondamentaux de Big Data -- 2025-2026}
\fancyfoot[C]{\thepage}
\renewcommand{\headrulewidth}{0.4pt}

% --- Encadres -----------------------------------------------------------
\usepackage[most]{tcolorbox}
\newtcolorbox{infobox}[1][]{
  colback=lightgray,
  colframe=sectiongreen,
  fonttitle=\bfseries\color{white},
  coltitle=white,
  colbacktitle=sectiongreen,
  arc=2mm, boxrule=0.6pt, breakable,
  left=8pt, right=8pt, top=6pt, bottom=6pt,
  #1}

% --- Acronymes -----------------------------------------------------------
\usepackage{acronym}

% =====================================================================
%                              DOCUMENT
% =====================================================================
\begin{document}

% --- Page de titre (a supprimer si le rapport est integre dans un autre fichier)
\begin{titlepage}
  \centering
  \vspace*{2cm}
  {\large\textbf{Université Abdelmalek Essaâdi}}\\
  {\large École Nationale des Sciences Appliquées -- Tétouan}\\
  {\normalsize Département : Intelligence Artificielle et Digitalisation}\\
  {\normalsize Filière : Sciences des Données, Big Data \& Intelligence Artificielle}\\
  {\normalsize Module : Fondamentaux de Big Data (M126)}\\
  {\normalsize Pr.~Imad Sassi}\\[2.5cm]
  {\Huge\bfseries\color{sectiongreen} Projet Final}\\[0.4cm]
  {\Large\bfseries\color{titlered} Fondamentaux de Big Data}\\[1.5cm]
  {\large\itshape Analyse prédictive de la charge réseau\\et optimisation de la consommation énergétique en temps réel}\\[3cm]
  {\large\textbf{Sections : 15, 17, 18 et 19}}\\[0.5cm]
  {\large Éthique \& Conformité --- Recommandations --- Perspectives --- Conclusion}\\[3cm]
  {\large\textbf{Année universitaire 2025--2026}}\\
  \vfill
\end{titlepage}

\tableofcontents
\newpage

% =====================================================================
% SECTION 15 : ETHIQUE ET CONFORMITE
% =====================================================================
\section{Éthique \& Conformité}
\label{sec:ethique}

\subsection{Introduction et enjeux}

Le projet manipule des données issues de plus de \textbf{91~000 compteurs intelligents} simulés à Tétouan, répartis sur \textbf{18 quartiers} et \textbf{681 distributeurs}. Bien que la simulation actuelle ne contienne aucune donnée nominative réelle, le système est conçu pour pouvoir, à terme, traiter de véritables relevés issus de compteurs Linky ou Smart~Meters. Cette section anticipe donc les obligations légales et éthiques applicables à un \emph{déploiement en production}.

L'énergie est un secteur sensible où les données fines de consommation peuvent révéler le \textbf{mode de vie} des habitants : heures de présence, nombre de personnes au foyer, équipements utilisés, périodes de vacances. La protection de ces données est donc une exigence forte, à la fois juridique et morale.

\begin{infobox}[title=Principe directeur du projet]
\textit{Aucune donnée ne doit pouvoir être reliée à un foyer ou à un individu identifiable. Toute information stockée à long terme est agrégée à l'échelle d'un quartier ou d'un distributeur.}
\end{infobox}

\subsection{Cadre juridique applicable}

\subsubsection{Au Maroc -- Loi 09-08}
La \textbf{Loi 09-08 (2009)} relative à la protection des personnes physiques à l'égard du traitement des données à caractère personnel est l'instrument juridique de référence. Elle est supervisée par la \textbf{CNDP} (Commission Nationale de contrôle de la Protection des Données à caractère Personnel). Elle impose notamment :
\begin{itemize}
  \item la \textbf{déclaration préalable} du traitement à la CNDP ;
  \item le \textbf{consentement explicite} des abonnés ;
  \item le \textbf{droit d'accès, de rectification et d'opposition} des personnes concernées ;
  \item la \textbf{limitation de la durée de conservation} des données ;
  \item la \textbf{sécurité et la confidentialité} des données traitées.
\end{itemize}

\subsubsection{Conformité internationale -- RGPD}
Si une partie des données transite par des serveurs européens (cloud, API météo telles que OpenWeather), le \textbf{RGPD} (Règlement UE~2016/679) s'applique de plein droit. Les principes structurants à retenir sont :
\begin{itemize}
  \item \textbf{Privacy by Design} : la confidentialité est intégrée dès la conception du système ;
  \item \textbf{Privacy by Default} : les paramètres les plus protecteurs sont activés par défaut ;
  \item \textbf{Minimisation des données} : seules les variables nécessaires à la finalité sont collectées ;
  \item \textbf{Pseudonymisation et chiffrement} systématiques ;
  \item Notification des violations de données dans un délai de \textbf{72~heures}.
\end{itemize}

\subsubsection{Norme sectorielle -- IEC 62056 (DLMS/COSEM)}
Standard international pour la communication des compteurs intelligents, qui impose des mécanismes d'authentification mutuelle et de chiffrement des trames échangées entre concentrateurs et compteurs.

\subsection{Mesures appliquées dans notre projet}

\subsubsection{Anonymisation et pseudonymisation}
Aucun identifiant nominatif n'est utilisé. Les compteurs sont identifiés par un code purement technique :

\begin{table}[h!]
\centering
\begin{tabular}{lll}
\toprule
\rowcolor{lightgray}
\textbf{Type d'entité} & \textbf{Format d'identifiant} & \textbf{Exemple} \\
\midrule
Compteur intelligent & \texttt{SM-TET-XXXXXX} & \texttt{SM-TET-045213} \\
Distributeur & \texttt{DIST-TET-XXXX} & \texttt{DIST-TET-0124} \\
Concentrateur & \texttt{CONC-TET-XX} & \texttt{CONC-TET-07} \\
\bottomrule
\end{tabular}
\caption{Format des identifiants techniques utilisés dans le projet}
\end{table}

Aucun nom, adresse, numéro de téléphone ou adresse e-mail n'est stocké dans MongoDB. La géolocalisation est effectuée \textbf{au niveau du quartier} (18 zones), jamais à l'échelle d'un foyer individuel.

\subsubsection{Minimisation des données}
Seules les variables strictement nécessaires à la prédiction de la charge sont collectées et stockées :
\begin{itemize}
  \item Mesures électriques : \texttt{energy\_consumption}, \texttt{voltage}, \texttt{current}, \texttt{status} ;
  \item Localisation grossière : \texttt{district}, \texttt{zone\_type}, latitude/longitude au niveau du quartier ;
  \item Horodatage : \texttt{timestamp}.
\end{itemize}
Aucune donnée comportementale fine (signature électrique d'appareils via \emph{Non-Intrusive Load Monitoring}) n'est extraite.

\subsubsection{Agrégation comme protection structurelle}
Le pipeline Spark Structured Streaming agrège les relevés par fenêtre de \textbf{15 minutes} et par quartier. Cette \textbf{double agrégation temporelle et spatiale} rend impossible la ré-identification d'un foyer individuel à partir des données stockées long terme dans la collection \texttt{spark\_district\_consumption}.

\subsubsection{Politique de rétention}
Le consumer Kafka est configuré avec la variable d'environnement \texttt{DELETE\_RAW\_AFTER\_COMPLETION="true"} : les relevés bruts sont \textbf{supprimés} de MongoDB après leur agrégation. Seules les agrégations par distributeur et par quartier sont conservées long terme. Cela répond directement au principe juridique de \textbf{limitation de la conservation}.

\subsubsection{Sécurité technique}

\begin{table}[h!]
\centering
\begin{tabularx}{\textwidth}{lYY}
\toprule
\rowcolor{lightgray}
\textbf{Couche} & \textbf{Mesure actuelle (dev)} & \textbf{Recommandation production} \\
\midrule
Réseau Docker & Réseau isolé \texttt{energy-net} & Politique réseau Kubernetes / VPN dédié \\
Kafka & PLAINTEXT (développement local) & TLS + SASL/SCRAM obligatoires \\
MongoDB & Sans authentification (dev) & Authentification + chiffrement TLS \\
Dashboard Streamlit & Accès ouvert local & OAuth2 / SSO + HTTPS \\
Logs applicatifs & \texttt{stdout} Docker & Centralisation ELK + masquage des PII \\
\bottomrule
\end{tabularx}
\caption{Comparaison entre les mesures de sécurité actuelles et celles requises en production}
\end{table}

\subsubsection{Transparence et explicabilité}
Le projet compare \textbf{trois modèles d'apprentissage automatique} (régression linéaire, Random Forest, Gradient Boosting) en exposant publiquement les métriques RMSE, MAE et R$^2$ via la collection MongoDB \texttt{ml\_model\_metrics}. Cette traçabilité répond directement au \textbf{droit à l'explication} prévu par l'article~22 du RGPD : un client doit pouvoir comprendre pourquoi le système prédit une saturation sur sa zone.

\subsection{Risques résiduels et plan de mitigation}

\begin{table}[h!]
\centering
\begin{tabularx}{\textwidth}{Y c Y}
\toprule
\rowcolor{lightgray}
\textbf{Risque identifié} & \textbf{Niveau} & \textbf{Mitigation} \\
\midrule
Ré-identification par croisement avec d'autres bases & Moyen & Agrégation à l'échelle quartier, k-anonymat $\geq 50$ \\
Biais géographique (quartiers riches mieux modélisés) & Moyen & Audit des métriques par zone, rééquilibrage \\
Fuite de données via API non chiffrée & Élevé en prod & TLS de bout en bout, audit régulier \\
Décisions automatisées discriminantes & Moyen & Validation humaine pour les alertes critiques \\
Dérive du modèle (data drift) & Moyen & Monitoring continu, ré-entraînement automatique \\
\bottomrule
\end{tabularx}
\caption{Cartographie des risques éthiques et des mesures de mitigation}
\end{table}

\subsection{Conclusion de la section}
Le projet respecte, par construction, les principes fondamentaux de la protection des données : minimisation, agrégation, pseudonymisation, durée de conservation limitée. Pour un passage en production réel sur le réseau de l'\textbf{ONEE} (Office National de l'Électricité et de l'Eau Potable), trois chantiers prioritaires subsistent :
\begin{enumerate}
  \item L'\textbf{authentification} systématique de tous les services internes (Kafka, MongoDB, Spark).
  \item Le \textbf{chiffrement TLS} de bout en bout pour tous les flux réseau.
  \item La \textbf{déclaration formelle} du traitement à la CNDP avant mise en service.
\end{enumerate}

\newpage

% =====================================================================
% SECTION 17 : RECOMMANDATIONS POUR LES PROFESSIONNELS
% =====================================================================
\section{Recommandations pour les professionnels}
\label{sec:recommandations}

\subsection{Introduction et publics cibles}

Le système livré ne se limite pas à un prototype académique : il fournit des leviers d'action concrets pour les acteurs du secteur énergétique marocain. Cette section formalise des recommandations adressées à \textbf{quatre publics professionnels} distincts, chacun pouvant exploiter différemment les sorties du pipeline (prédictions, alertes, corrélations NLP, métriques de qualité) :

\begin{enumerate}
  \item Les \textbf{gestionnaires de réseau de distribution} (ONEE, Amendis, Lydec, Redal) ;
  \item Les \textbf{fournisseurs et traders d'énergie} ;
  \item Les \textbf{équipes IT et Data Science} en charge de l'exploitation du système ;
  \item Les \textbf{décideurs publics et régulateurs} (Ministère de la Transition Énergétique, CNDP, ANRE).
\end{enumerate}

Chaque recommandation est justifiée par les comportements observés dans le projet et par les sorties effectivement produites par le pipeline (collections MongoDB \texttt{spark\_district\_consumption}, \texttt{ml\_predictions}, \texttt{alerts}, \texttt{nlp\_correlations}).

\subsection{Recommandations pour les gestionnaires de réseau}

\subsubsection{Pilotage opérationnel temps réel}
Le dashboard Streamlit permet de surveiller en continu les 18 quartiers de Tétouan. Nous recommandons aux opérateurs de :
\begin{itemize}
  \item Mettre en place une \textbf{salle de contrôle 24/7} avec rotation des équipes selon le flux d'alertes prédit.
  \item Définir des \textbf{seuils d'action gradués} (vert / orange / rouge) intégrés directement dans le dashboard.
  \item Documenter une \textbf{procédure d'escalade} pour chaque type d'alerte (\texttt{voltage\_risk}, \texttt{overload\_risk}, \texttt{saturation\_risk}).
  \item Conserver un \textbf{journal d'incidents} alimenté automatiquement par les alertes pour audit a posteriori.
\end{itemize}

\begin{table}[h!]
\centering
\begin{tabularx}{\textwidth}{lcccY}
\toprule
\rowcolor{lightgray}
\textbf{KPI} & \textbf{Vert} & \textbf{Orange} & \textbf{Rouge} & \textbf{Action recommandée} \\
\midrule
Tension moyenne (V) & $>$ 220 & 210--220 & $<$ 210 & Investigation transformateur immédiate \\
Saturation zone (\%) & $<$ 60 & 60--80 & $>$ 80 & Préparer délestage proactif \\
Compteurs en panne (\%) & $<$ 1 & 1--5 & $>$ 5 & Dépêche équipe technique \\
Anomalies NLP / 15 min & $<$ 5 & 5--15 & $>$ 15 & Renforcer hotline support \\
Écart prévu vs réel (\%) & $<$ 5 & 5--15 & $>$ 15 & Revue du modèle ML \\
\bottomrule
\end{tabularx}
\caption{Seuils opérationnels recommandés pour la salle de contrôle}
\end{table}

\subsubsection{Délestage proactif intelligent}
Plutôt que de subir les coupures, l'horizon de prédiction de 15~minutes permet d'organiser un \textbf{délestage sélectif et hiérarchisé} :
\begin{itemize}
  \item Établir des \textbf{contrats d'effacement} avec les gros clients industriels en échange de tarifs préférentiels.
  \item Définir un \textbf{ordre de priorité} : industries volontaires $\rightarrow$ commerces $\rightarrow$ tertiaire $\rightarrow$ résidentiel (dernier recours).
  \item Activer le délestage \textbf{avant la saturation effective}, lorsque la prédiction franchit le seuil orange.
  \item Communiquer en \textbf{temps réel} aux clients délestés (SMS, application mobile).
\end{itemize}

\subsubsection{Maintenance prédictive des transformateurs}
La corrélation NLP/tension implémentée dans le projet (\texttt{nlp\_correlations}) est un \textbf{indicateur précoce} de défaillance d'équipement :
\begin{itemize}
  \item Toute zone présentant simultanément une \textbf{dérive de tension} ($<$ 215~V) et un \textbf{pic de mots-clés négatifs} (\emph{panne, coupure, surcharge}) doit déclencher une inspection physique.
  \item Programmer la maintenance préventive pendant les \textbf{heures de faible charge} (typiquement 02:00--05:00) identifiées dans les agrégations Spark.
  \item Tenir un \textbf{historique pluriannuel} pour entraîner un futur modèle de \emph{Time-To-Failure}.
\end{itemize}

\subsection{Recommandations pour les fournisseurs d'énergie}

\subsubsection{Optimisation des achats sur le marché de gros}
Les prédictions à 15~minutes peuvent être agrégées pour produire des \textbf{prévisions \emph{day-ahead}} à l'échelle régionale :
\begin{itemize}
  \item Affiner les \textbf{nominations} sur le marché de gros (réduction des pénalités de déséquilibre).
  \item Ajuster les \textbf{stratégies de couverture} (\emph{hedging}) selon les prévisions météo intégrées.
  \item Identifier les \textbf{créneaux d'arbitrage} entre production locale et import.
\end{itemize}

\subsubsection{Tarification dynamique (Demand Response)}
Recommandations pour évoluer du modèle binaire heures-pleines/creuses vers un modèle dynamique :
\begin{itemize}
  \item Définir des \textbf{tranches tarifaires variables} basées sur la charge prédite par quartier.
  \item Notifier les clients via \textbf{application mobile} et SMS lorsque le tarif passe en heures creuses dynamiques.
  \item Introduire des \textbf{mécanismes de récompense} (cashback, points fidélité) pour les clients qui décalent leur consommation.
  \item Cibler en priorité les usages \textbf{décalables} : lave-linge, ballon d'eau chaude, recharge de véhicules électriques.
\end{itemize}

\subsubsection{Transparence et relation client}
\begin{itemize}
  \item Exploiter les agrégations pour offrir une \textbf{facture détaillée} avec décomposition horaire.
  \item Fournir des \textbf{recommandations personnalisées} d'économie d'énergie (« \emph{vous consommez 30~\% de plus que la moyenne de votre quartier} »).
  \item Garantir le \textbf{droit à l'explication} : tout client doit pouvoir comprendre pourquoi son quartier a été délesté ou facturé à un tarif majoré.
\end{itemize}

\subsection{Recommandations pour les équipes IT et Data Science}

\subsubsection{Exploitation et MLOps}
Pour assurer la pérennité du système au-delà du prototype :
\begin{itemize}
  \item Mettre en place un \textbf{ré-entraînement périodique} des trois modèles ML (mensuel au minimum, hebdomadaire si dérive détectée).
  \item Monitorer la \textbf{dérive conceptuelle} (\emph{concept drift}) via la collection \texttt{ml\_model\_metrics} et alerter si la R$^2$ chute de plus de 10~\%.
  \item Industrialiser le \textbf{registre de modèles} (MLflow) pour tracer chaque version déployée.
  \item Automatiser les \textbf{sauvegardes MongoDB} (snapshots quotidiens, rétention 30~jours).
  \item Garantir la \textbf{haute disponibilité Kafka} via un cluster 3 brokers en production.
\end{itemize}

\subsubsection{Sécurité et conformité opérationnelle}
\begin{itemize}
  \item Activer \textbf{TLS} sur tous les flux Kafka et MongoDB dès la mise en production.
  \item Mettre en place une \textbf{authentification SASL/SCRAM} entre producers et brokers.
  \item Effectuer des \textbf{audits de sécurité} semestriels (tests d'intrusion, revue de code).
  \item Désigner un \textbf{Délégué à la Protection des Données} (DPO) en interne.
  \item Documenter un \textbf{plan de réponse aux incidents} conforme au délai de notification de 72~heures.
\end{itemize}

\subsubsection{Documentation et transmission de connaissance}
\begin{itemize}
  \item Maintenir des \textbf{runbooks opérationnels} pour chaque incident type (broker Kafka down, lag de consumer, pic de latence Spark).
  \item Documenter l'\textbf{onboarding} des nouveaux ingénieurs sur la stack.
  \item Organiser des \textbf{revues post-mortem} systématiques après chaque incident majeur.
  \item Publier en interne les \textbf{indicateurs de santé} du système (SLA, MTTR, MTBF).
\end{itemize}

\subsection{Recommandations pour les décideurs publics et régulateurs}

\subsubsection{Cadre réglementaire favorable aux smart grids}
\begin{itemize}
  \item Standardiser les \textbf{formats d'échange} de données entre opérateurs (DLMS/COSEM, IEC~61850).
  \item Imposer un \textbf{déploiement progressif} des compteurs intelligents sur tout le territoire (à l'image du programme Linky en France).
  \item Encadrer la \textbf{portabilité des données} pour les consommateurs (équivalent du droit RGPD).
  \item Faciliter le \textbf{partage de données anonymisées} entre opérateurs pour améliorer les modèles à l'échelle nationale.
\end{itemize}

\subsubsection{Politique tarifaire et incitations}
\begin{itemize}
  \item Mettre en place une \textbf{fiscalité incitative} pour les équipements communicants (chauffe-eau pilotables, bornes de recharge intelligentes).
  \item Subventionner l'\textbf{autoconsommation collective} dans les quartiers à forte concentration de toitures solaires.
  \item Encadrer les \textbf{contrats Demand Response} pour protéger les clients vulnérables.
\end{itemize}

\subsubsection{Soutien à l'objectif national 52~\% EnR en 2030}
\begin{itemize}
  \item Financer le \textbf{stockage stationnaire} (batteries lithium, STEP) indispensable pour absorber l'intermittence.
  \item Accélérer le \textbf{raccordement} des projets EnR au réseau de distribution.
  \item Encourager la \textbf{recherche académique} (modèles prédictifs, jumeaux numériques du réseau).
\end{itemize}

\subsubsection{Protection du consommateur}
\begin{itemize}
  \item Renforcer les contrôles de la \textbf{CNDP} sur l'usage des données de consommation.
  \item Imposer un \textbf{droit à l'explication} pour toute décision automatisée (délestage, tarif majoré).
  \item Garantir un \textbf{tarif social} protégé contre la tarification dynamique pour les ménages précaires.
\end{itemize}

\subsection{Tableau de bord d'indicateurs (KPI) recommandés}

Pour piloter durablement le système, nous proposons un tableau de bord articulé autour de quatre familles d'indicateurs :

\begin{table}[h!]
\centering
\begin{tabularx}{\textwidth}{lYY}
\toprule
\rowcolor{lightgray}
\textbf{Famille} & \textbf{Indicateur} & \textbf{Cible recommandée} \\
\midrule
\multirow{3}{*}{\textbf{Opérationnel}}
  & Disponibilité du dashboard & $\geq$ 99,5~\% \\
  & Latence end-to-end (compteur $\rightarrow$ alerte) & $<$ 30~secondes \\
  & Taux de fausses alertes & $<$ 5~\% \\
\midrule
\multirow{3}{*}{\textbf{Métier}}
  & Réduction des pannes par surcharge & $-$30~\% à 12~mois \\
  & Taux d'effacement réussi (Demand Response) & $\geq$ 70~\% \\
  & Taux de pénétration EnR par quartier & + 5~\% / an \\
\midrule
\multirow{3}{*}{\textbf{Modèles ML}}
  & RMSE de prédiction & $<$ 8~\% de la charge moyenne \\
  & R$^2$ des trois modèles & $\geq$ 0,85 \\
  & Délai de ré-entraînement & $<$ 30~jours \\
\midrule
\multirow{3}{*}{\textbf{Éthique \& sécurité}}
  & Incidents de sécurité (par an) & 0 incident majeur \\
  & Délai de notification CNDP & $<$ 72~heures \\
  & Taux d'agrégation des données stockées & 100~\% \\
\bottomrule
\end{tabularx}
\caption{Tableau de bord des KPIs recommandés pour le pilotage du système}
\end{table}

\subsection{Synthèse et priorisation des actions}

Pour faciliter l'appropriation des recommandations par les équipes opérationnelles, nous proposons une priorisation à trois horizons :

\begin{table}[h!]
\centering
\begin{tabularx}{\textwidth}{lY}
\toprule
\rowcolor{lightgray}
\textbf{Horizon} & \textbf{Actions prioritaires} \\
\midrule
\textbf{Court terme} (0--3 mois) &
  Déployer le dashboard en salle de contrôle. Activer TLS sur Kafka/Mongo. Définir les seuils d'alerte. Désigner le DPO. \\
\midrule
\textbf{Moyen terme} (3--12 mois) &
  Industrialiser MLOps (MLflow + ré-entraînement automatique). Lancer un pilote de tarification dynamique sur 1~quartier. Mettre en place les contrats d'effacement industriels. \\
\midrule
\textbf{Long terme} ($>$ 12 mois) &
  Étendre à l'échelle régionale puis nationale. Intégrer les données EnR (PVGIS, éolien). Évoluer vers un véritable Smart Grid avec pilotage automatisé. \\
\bottomrule
\end{tabularx}
\caption{Plan d'action priorisé pour les acteurs du secteur énergétique}
\end{table}

\begin{infobox}[title=Message clé]
\textit{Les bénéfices du projet ne se concrétisent que si les recommandations ci-dessus sont mises en \oe{}uvre conjointement par les opérateurs, les régulateurs et les équipes techniques. La technologie seule ne suffit pas : c'est la \textbf{combinaison} d'une infrastructure Big Data robuste, d'une gouvernance éthique des données et d'une volonté politique forte qui permettra au Maroc de réussir sa transition énergétique.}
\end{infobox}

\newpage

% =====================================================================
% SECTION 18 : PERSPECTIVES D'AMELIORATION
% =====================================================================
\section{Perspectives d'amélioration}
\label{sec:perspectives}

\subsection{Limites identifiées du projet actuel}

Avant de proposer des évolutions, il est essentiel d'identifier honnêtement les limites du système livré :

\begin{table}[h!]
\centering
\begin{tabularx}{\textwidth}{lY}
\toprule
\rowcolor{lightgray}
\textbf{Limite} & \textbf{Description} \\
\midrule
Données simulées & Le simulateur génère des consommations selon des profils théoriques. Les comportements réels présentent davantage de variabilité. \\
Volumétrie testée & 91~000 compteurs en mode pseudo-distributed Docker. Le passage à un cluster physique multi-n\oe{}uds n'a pas été éprouvé. \\
Modèles ML classiques & scikit-learn : LinearRegression / RandomForest / GradientBoosting. Pas de modèles séquentiels (LSTM, Transformer). \\
Fenêtre fixe de 15~minutes & Pas d'analyse fine intra-fenêtre, ni de granularité variable. \\
NLP limité au français & TF-IDF + Logistic Regression, sans embeddings contextuels (BERT) ni gestion de l'arabe dialectal. \\
Pas d'intégration EnR & Les énergies renouvelables (solaire, éolien) ne sont pas encore introduites comme features prédictives. \\
Pas de boucle de rétroaction & Le système prédit mais ne pilote pas encore d'actionneurs (délestage automatique, tarification dynamique). \\
\bottomrule
\end{tabularx}
\caption{Limites assumées du projet dans sa version actuelle}
\end{table}

\subsection{Axes d'amélioration techniques}

\subsubsection{Passage à des données réelles}
\begin{itemize}
  \item Établir un \textbf{partenariat avec l'ONEE} pour obtenir un échantillon anonymisé de relevés Linky/Smart~Meters réels sur Tétouan.
  \item Mise en place d'un protocole \textbf{DLMS/COSEM (IEC 62056)} pour la collecte directe depuis les concentrateurs physiques.
  \item Comparaison statistique des distributions simulées vs réelles pour calibrer le simulateur et réduire le \emph{simulation-to-reality gap}.
\end{itemize}

\subsubsection{Modèles d'IA plus puissants}

\begin{table}[h!]
\centering
\begin{tabularx}{\textwidth}{lYY}
\toprule
\rowcolor{lightgray}
\textbf{Catégorie} & \textbf{Modèle proposé} & \textbf{Apport attendu} \\
\midrule
Séries temporelles & LSTM / GRU & Capturer les dépendances long terme (cycles hebdomadaires, saisonniers) \\
Mécanismes d'attention & Transformer (Informer, Autoformer) & Prédiction multi-horizon précise \\
Modèles hybrides & N-BEATS, Prophet + XGBoost & Combinaison de tendance + événements \\
Apprentissage fédéré & TensorFlow Federated & Entraînement sans centraliser les données (vie privée) \\
Détection d'anomalies & Isolation Forest, Autoencoders & Identifier des incidents jamais vus \\
\bottomrule
\end{tabularx}
\caption{Familles de modèles avancés envisagés}
\end{table}

\subsubsection{Enrichissement du pipeline NLP}
\begin{itemize}
  \item Remplacement de TF-IDF par \textbf{CamemBERT} ou \textbf{DarijaBERT} pour traiter les feedbacks en français standard et en arabe dialectal marocain.
  \item Ajout d'une couche d'\textbf{analyse de sentiment} (détresse, colère) pour prioriser les incidents critiques.
  \item \textbf{Reconnaissance d'entités nommées} (NER) pour relier automatiquement un rapport à un quartier, un transformateur ou un distributeur précis.
  \item Construction d'un \textbf{graphe de connaissances} (Neo4j) reliant incidents, équipements et zones géographiques.
\end{itemize}

\subsubsection{Architecture et scalabilité}
\begin{itemize}
  \item Migration de Docker~Compose vers \textbf{Kubernetes} (K3s ou OpenShift) pour le mode \emph{Fully-Distributed} à grande échelle.
  \item Remplacement du single-broker Kafka par un \textbf{cluster Kafka 3 brokers} avec \texttt{replication.factor = 3}.
  \item Adoption d'\textbf{Apache Flink} en complément de Spark pour les latences sub-secondes (alertes critiques temps réel).
  \item Stockage analytique : ajout d'\textbf{Apache Druid} ou \textbf{ClickHouse} pour les requêtes OLAP du dashboard.
  \item Data lake : intégration d'\textbf{Apache Iceberg} sur S3/MinIO pour l'archivage long terme.
\end{itemize}

\subsubsection{Qualité des données et MLOps}
\begin{itemize}
  \item Mise en place de \textbf{Great Expectations} pour valider chaque flux entrant avant ingestion.
  \item \textbf{Feature Store} (Feast) pour partager les variables prédictives entre équipes et entre modèles.
  \item \textbf{MLflow} pour le suivi des expérimentations et la gestion d'un \emph{model registry}.
  \item Ré-entraînement automatique des modèles via \textbf{Apache Airflow} dès que la dérive (\emph{concept drift}) dépasse un seuil prédéfini.
  \item Tests unitaires et d'intégration dans une chaîne CI/CD GitHub~Actions.
\end{itemize}

\subsection{Axes d'amélioration métier}

\subsubsection{Intégration des Énergies Renouvelables (EnR)}
\begin{itemize}
  \item Ajout d'un producer \textbf{solaire} consommant l'API \emph{PVGIS} de la JRC européenne, paramétrée pour Tétouan (latitude 35.57°N).
  \item Ajout d'un producer \textbf{éolien} combinant vitesse du vent et courbes de puissance des turbines.
  \item Modèle de \textbf{prévision nette} : consommation prévue $-$ production EnR prévue.
  \item Recommandation d'\textbf{heures creuses dynamiques} pour favoriser l'autoconsommation.
\end{itemize}

\subsubsection{Pilotage actif du réseau (Smart Grid)}
\begin{itemize}
  \item \textbf{Délestage intelligent} : envoi automatique de commandes vers les distributeurs en cas de saturation prévue.
  \item \textbf{Tarification dynamique} : ajustement du prix au quart d'heure selon la charge prédite (\emph{Demand Response}).
  \item \textbf{Recommandations personnalisées} aux clients via une application mobile : par exemple « Décalez votre lave-linge de 20:00 à 23:00 pour économiser X~MAD ».
  \item Intégration avec les \textbf{bornes de recharge de véhicules électriques} pour lisser les pics de consommation.
\end{itemize}

\subsubsection{Extension géographique}
\begin{itemize}
  \item Réplication du système sur d'autres villes marocaines : \textbf{Tanger, Casablanca, Rabat, Marrakech}.
  \item Couche d'agrégation \textbf{régionale} puis \textbf{nationale} pour piloter le réseau ONEE complet.
  \item Intégration avec les \textbf{interconnexions internationales} (Espagne, Algérie).
\end{itemize}

\subsubsection{Nouveaux cas d'usage analytiques}
\begin{itemize}
  \item \textbf{Détection de fraude} : identification de schémas de consommation anormaux pouvant révéler un vol d'électricité.
  \item \textbf{Maintenance prédictive} des transformateurs à partir des dérives de tension observées.
  \item \textbf{Étude d'impact climatique} : corrélation entre vagues de chaleur et pics de consommation.
  \item \textbf{Empreinte carbone} par quartier en croisant la consommation avec le mix énergétique en temps réel.
\end{itemize}

\subsection{Axes d'amélioration de l'interface}

\begin{table}[h!]
\centering
\begin{tabularx}{\textwidth}{lY}
\toprule
\rowcolor{lightgray}
\textbf{Évolution} & \textbf{Description} \\
\midrule
Application mobile & React Native pour les techniciens terrain (alertes \emph{push}) \\
Mode sombre + accessibilité & Conformité WCAG 2.1 niveau AA \\
Multi-langue & Français, arabe, anglais, tamazight \\
Export PDF automatique & Rapport hebdomadaire généré pour les gestionnaires de réseau \\
Chatbot LLM & Interrogation du dashboard en langage naturel («~Quelles zones risquent une saturation ce soir~?~») \\
\bottomrule
\end{tabularx}
\caption{Évolutions ergonomiques envisagées pour le dashboard}
\end{table}

\subsection{Feuille de route proposée}

\begin{table}[h!]
\centering
\begin{tabularx}{\textwidth}{lYY}
\toprule
\rowcolor{lightgray}
\textbf{Horizon} & \textbf{Priorité 1} & \textbf{Priorité 2} \\
\midrule
0--3 mois & Sécurisation TLS Kafka + MongoDB & Intégration BERT pour le NLP \\
3--6 mois & Ajout LSTM et comparaison aux 3 modèles actuels & Producer solaire (PVGIS) \\
6--12 mois & Migration Kubernetes + cluster multi-n\oe{}uds & Application mobile React Native \\
12 mois et + & Apprentissage fédéré + déploiement réel ONEE & Extension à 5 villes marocaines \\
\bottomrule
\end{tabularx}
\caption{Roadmap d'évolution du système sur 12 mois et au-delà}
\end{table}

\newpage

% =====================================================================
% SECTION 19 : CONCLUSION
% =====================================================================
\section{Conclusion}
\label{sec:conclusion}

\subsection{Bilan du projet}

Ce projet a permis de \textbf{concevoir, implémenter et valider} un pipeline Big Data complet répondant aux \textbf{huit exigences techniques} du sujet (points A à H) :

\begin{itemize}[leftmargin=*]
  \item[\done] \textbf{Ingestion en streaming} de trois sources hétérogènes (compteurs, météo, RSS industriel) via Apache Kafka.
  \item[\done] \textbf{Traitement distribué} par Spark Structured Streaming avec fenêtrage de 15~minutes et corrélation NLP/tension.
  \item[\done] \textbf{Stockage NoSQL} dans MongoDB, choisi à l'issue d'une étude comparative avec Cassandra, HBase, Neo4j et ElasticSearch.
  \item[\done] \textbf{Modélisation IA} comparant rigoureusement trois algorithmes (régression linéaire, Random Forest, Gradient Boosting) sur RMSE, MAE, R$^2$ et temps d'inférence.
  \item[\done] \textbf{Data Quality Framework} couvrant complétude, bruit et biais, avec choix assumé de données synthétiques.
  \item[\done] \textbf{Dashboard interactif} Streamlit affichant carte de chaleur, courbes prédiction~vs~réel, nuage de mots et alertes prédictives.
  \item[\done] \textbf{Modes d'installation hybride} testés en \emph{Pseudo-Distributed} (Docker~Compose) et \emph{Fully-Distributed} (Spark master + workers).
  \item[\done] \textbf{Éthique et conformité} anticipées (anonymisation, agrégation, RGPD, Loi 09-08).
\end{itemize}

\subsection{Réponse à la problématique}

La problématique posée était la suivante :
\begin{quote}
\itshape Comment prédire avec précision la charge réseau et détecter en temps réel les incidents pour optimiser la distribution énergétique et favoriser l'intégration des énergies renouvelables~?
\end{quote}

Le système livré apporte une réponse concrète et opérationnelle :
\begin{itemize}
  \item Prédiction par quartier à un horizon de \textbf{15 minutes}, avec une RMSE compétitive sur les trois modèles évalués.
  \item Détection en temps réel des \textbf{chutes de tension} et des \textbf{surcharges} via le simulateur Tétouan.
  \item Corrélation automatique entre \textbf{anomalies textuelles} (rapports techniques, flux RSS) et \textbf{anomalies physiques} (\texttt{avg\_voltage} $<$ 210~V), permettant un diagnostic plus précis qu'une approche purement quantitative.
  \item Génération d'\textbf{alertes prédictives} sur les zones à risque (\texttt{saturation\_risk}, \texttt{voltage\_risk}, \texttt{overload\_risk}).
\end{itemize}

\subsection{Apports du projet}

\subsubsection{Apports techniques}
\begin{itemize}
  \item Maîtrise d'une \textbf{stack Big Data complète} : Kafka, Spark, MongoDB, Docker, Streamlit.
  \item Mise en pratique des paradigmes \textbf{streaming temps réel} et \textbf{batch analytics}.
  \item Comparaison rigoureuse de modèles ML accompagnée d'une \textbf{étude de performances}.
  \item Intégration de techniques de \textbf{NLP en français} dans une chaîne décisionnelle métier.
\end{itemize}

\subsubsection{Apports métier}
\begin{itemize}
  \item Outillage utilisable par un \textbf{gestionnaire de réseau} pour anticiper la demande énergétique.
  \item Recommandations exploitables pour :
    \begin{itemize}
      \item le \textbf{délestage proactif} des zones à risque ;
      \item l'\textbf{ajustement tarifaire} selon la charge prévue ;
      \item la \textbf{planification de la maintenance} des transformateurs ;
      \item l'\textbf{intégration des EnR intermittentes} dans le mix énergétique national.
    \end{itemize}
\end{itemize}

\subsubsection{Apports académiques}
\begin{itemize}
  \item Application concrète des principes étudiés dans le module M126 (\emph{Fondamentaux de Big Data}).
  \item Démonstration de la complémentarité entre \textbf{streaming distribué} et \textbf{machine learning}.
  \item Illustration d'un cas réel d'\textbf{éthique du Big Data} dans un secteur fortement régulé.
\end{itemize}

\subsection{Limites assumées}

\begin{itemize}
  \item Le projet repose sur des \textbf{données simulées} et non sur des relevés Linky réels, faute d'accès aux infrastructures ONEE.
  \item La validation des modèles s'appuie sur un échantillon limité de cycles ; le système nécessite plusieurs jours d'historique avant d'atteindre sa pleine performance prédictive.
  \item Le déploiement \emph{Fully-Distributed} a été simulé via Docker~Compose, sans cluster physique multi-n\oe{}uds dédié.
  \item Les modèles utilisés restent \textbf{classiques} (scikit-learn) ; les approches deep learning sont laissées en perspective.
\end{itemize}

\subsection{Mot de fin}

Au-delà de la note académique, ce projet constitue une \textbf{base technique réutilisable} pour des cas d'usage réels d'optimisation des réseaux énergétiques au Maroc. À l'heure où le Royaume vise \textbf{52~\% d'énergies renouvelables d'ici 2030}, des systèmes capables de prévoir la charge et d'orchestrer un réseau intelligent deviennent stratégiques pour la transition énergétique nationale.

Le code source, les modèles entraînés, le dashboard et la présente documentation sont publiés sur le dépôt \textbf{GitHub} du groupe, dans une logique de \textbf{reproductibilité} et de transparence scientifique.

\bigskip

\begin{infobox}[title=Remerciements]
Nous remercions chaleureusement \textbf{Pr.~Imad Sassi} pour l'encadrement pédagogique du module \emph{Fondamentaux de Big Data} et pour avoir proposé un sujet à la fois exigeant techniquement et porteur d'impact sociétal. Nos remerciements s'adressent également à l'\textbf{équipe pédagogique} du Département IA \& Digitalisation de l'\textbf{ENSA-Tétouan} pour la qualité de la formation dispensée.
\end{infobox}

\end{document}
